\subsection{Primitivity and rank-one trace matrices}

\begin{definition}[Explicit $\eta$-data from the sector tensors]%
    \label{def:mpdo_explicit_eta_of_sector_tensors}
    \lean{MPOTensor.ExplicitEtaOperators.ofSectorTensors}
    \leanok
    \uses{def:mpdo_sector_eta, def:mpdo_explicit_eta_operators}
    If every neighboring operator $\eta_{k,h}$ of
    Definition~\ref{def:mpdo_sector_eta} is positive semidefinite, the
    family constitutes explicit neighboring $\eta$-data. The source obtains
    the positivity from the projected chain identity: writing
    $\tilde\sigma
    = U^{\dagger\otimes N} \sigma^{(N)}({\cal K})\,U^{\otimes N}$
    for the $N$-site operator in the sector basis and $Q_k$ for the
    projector onto the $k$-th sector of the physical splitting, every
    cyclic sector assignment $k_1,\ldots,k_N$ (with $k_{N+1}=k_1$) satisfies
    \begin{align}
        0
        &\leq
        [Q_{k_1}\otimes\cdots\otimes Q_{k_N}]
        \tilde\sigma
        [Q_{k_1}\otimes\cdots\otimes Q_{k_N}]
        =\bigotimes_{n=1}^{N}\eta_{k_n,k_{n+1}},
        \notag
    \end{align}
    and the sector tensors can be chosen so that each factor is positive
    semidefinite \cite[Appendix~C.2, lines~1446--1450]{Cirac2016MPDO_arXiv}.
    The all-length sector identity is
    Theorem~\ref{thm:mpdo_sector_adapted_decomposition}. That algebraic
    identity does not by itself choose coherent positive representatives for
    all neighboring pairs, so positive semidefiniteness remains a hypothesis
    on the chosen representatives here.
    Theorem~\ref{thm:mpdo_sector_eta_pair_posSemidef} provides the diagonal
    and paired positivity, and
    Theorem~\ref{thm:mpdo_eta_positivity_choice_symmetric} a rescaled
    positive family under symmetric support.
\end{definition}

\begin{theorem}[Trace matrix of the sector-tensor $\eta$-data]
    \label{thm:mpdo_sector_eta_data_trace_matrix}
    \lean{MPOTensor.ExplicitEtaOperators.traceMatrix_ofSectorTensors}
    \leanok
    \uses{def:mpdo_explicit_eta_of_sector_tensors,
        def:mpdo_explicit_eta_trace_matrix, def:mpdo_closed_sector_tensors}
    Under the hypotheses of
    Definition~\ref{def:mpdo_explicit_eta_of_sector_tensors}, the trace
    matrix of the explicit $\eta$-data is the pairing of the closed sector
    tensors:
    \begin{align}
        T_{k,h}
        &=\tr(\eta_{k,h})=(r_k|l_h).
        \notag
    \end{align}
\end{theorem}

\begin{proof}
    \leanok
    \uses{thm:mpdo_sector_trace_pairing}
    Immediate from Theorem~\ref{thm:mpdo_sector_trace_pairing} and the
    definition of the trace matrix.
\end{proof}

\begin{theorem}[Primitivity from returns of lengths two and three]%
    \label{thm:matrix_primitive_of_two_three_returns}
    \lean{Matrix.IsIrreducible.%
      isPrimitive_of_pow_two_diagonal_pos_of_pow_three_diagonal_pos}
    \leanok
    Let $T$ be a non-negative irreducible matrix on a finite index set. If
    $(T^2)_{k,k}>0$ and $(T^3)_{k,k}>0$ for every index $k$, then $T$ is
    primitive.
\end{theorem}

\begin{proof}
    \leanok
    Choose a positive path between every ordered pair of indices, and let $L$
    be the sum of the chosen path lengths. Every chosen length is at most $L$.
    For a chosen path of length $m$, write the difference $2L+2-m$ as
    $2a+3b$ with $a,b\geq0$. Appending these returns at the terminal index
    gives a positive path of common length $2L+2$ between every ordered pair.
\end{proof}

\begin{definition}[Nonzero-weight sector trace matrix]%
    \label{def:mpdo_active_sector_trace_matrix}
    \lean{MPOTensor.PhysicalSectorFactorization.ActiveSector,
      MPOTensor.PhysicalSectorFactorization.ActiveSectorEntryIndex}
    \lean{MPOTensor.PhysicalSectorFactorization.activeSectorOneSiteMatrixFamily,
      MPOTensor.PhysicalSectorFactorization.activeSectorSet,
      MPOTensor.PhysicalSectorFactorization.activeSectorOneSiteSpace}
    \lean{MPOTensor.PhysicalSectorFactorization.activeSectorTraceMatrix}
    \leanok
    \uses{def:mpdo_one_site_sector_matrix_space,
      def:mpdo_minimal_neighboring_operator}
    For sector weights $p_k$, let $I_*=\{k:p_k\ne0\}$. The one-site matrices
    and their spaces on $I_*$ are
    \begin{align}
        A_{k;x,y},
        \notag\\
        {\cal A}_S
        &=\spn\{A_{k;x,y}:k\in S\},
        \qquad S\subseteq I_*.
        \notag
    \end{align}
    The real trace matrix on this index set is
    \begin{align}
        T^*_{k,h}
        &=\re\tr(\eta_{k,h}),
        \qquad k,h\in I_*.
        \notag
    \end{align}
    Only the trace matrix is restricted: the physical direct sum still
    contains every zero-weight summand.
\end{definition}

\begin{theorem}[Identification and spanning of active one-site matrices]%
    \label{thm:mpdo_active_one_site_matrix_spanning}
    \lean{MPOTensor.PhysicalSectorFactorization.%
      oneSiteSectorMatrix_eq_sectorVirtualMatrix}
    \lean{MPOTensor.PhysicalSectorFactorization.%
      activeSectorOneSiteMatrixFamily_span_eq_top}
    \leanok
    \uses{def:mpdo_active_sector_trace_matrix,
      thm:mpdo_sector_virtual_matrix_spanning}
    For every sector entry, the one-site matrix used in the directed-cut
    argument equals the virtual matrix used in the spanning argument. If all
    sector virtual matrices span $\MN{D}$ and the matrices in every
    zero-weight sector vanish, then
    $\spn\{A_{k;x,y}:k\in I_*\}=\MN{D}$.
\end{theorem}

\begin{theorem}[Complementary active-sector spaces]%
    \label{thm:mpdo_active_sector_complementary_spaces}
    \lean{MPOTensor.PhysicalSectorFactorization.%
      activeSectorOneSiteSpace_sup_compl_eq_top}
    \lean{MPOTensor.PhysicalSectorFactorization.activeSectorOneSiteSpace_ne_bot}
    \lean{MPOTensor.PhysicalSectorFactorization.%
      mul_eq_zero_of_mem_activeSectorOneSiteSpace}
    \leanok
    \uses{def:mpdo_active_sector_trace_matrix,
      thm:mpdo_active_one_site_matrix_spanning,
      thm:mpdo_directed_cut_one_site_spaces}
    Suppose that the active one-site matrices span $\MN{D}$ and each active
    sector contains a nonzero such matrix. For every $S\subseteq I_*$,
    ${\cal A}_S+{\cal A}_{I_*\setminus S}=\MN{D}$, and both spaces are
    nonzero when the corresponding sector sets are nonempty. If
    $\eta_{k,h}=0$ for $k\in S$ and $h\in C$, then
    ${\cal A}_S{\cal A}_C=0$.
\end{theorem}

\begin{theorem}[Strong connectivity of the active neighboring support]%
    \label{thm:mpdo_active_neighboring_support_strongly_connected}
    \lean{MPOTensor.PhysicalSectorFactorization.%
      activeSector_reflTransGen_neighboringOperator_ne_zero}
    \lean{MPOTensor.PhysicalSectorFactorization.%
      exists_active_neighboringOperator_ne_zero}
    \leanok
    \uses{thm:mpdo_active_sector_complementary_spaces,
      lem:matrix_one_sided_product_obstruction}
    Suppose that the active one-site matrices span $\MN{D}$ and each active
    sector contains a nonzero such matrix. Then for every $k,h\in I_*$ there
    is a directed walk
    \begin{align}
        k=k_0\longrightarrow k_1\longrightarrow\cdots
        \longrightarrow k_m=h,
        \notag
    \end{align}
    such that $\eta_{k_j,k_{j+1}}\ne0$. In particular, every active sector
    has an outgoing nonzero neighboring operator.
\end{theorem}

\begin{figure}[ht]
    \centering
    \begin{align}
        P_j\mathcal K_iP_{j^\prime}
        &=0\quad\text{unless }i=j=j^\prime.
        \notag
    \end{align}
    \begin{center}
        % P_j K_i P_{j'} = 0 unless i=j=j': the sector-projector sandwich of a
        % BNT element K_i.  Three `op' rows compose vertically: P_j's own down
        % leg pairs with K_i's up leg, and K_i's down leg pairs with P_{j'}'s
        % up leg, so only P_j's north stub and P_{j'}'s south stub survive --
        % the matrix-product indices of P_j K_i P_{j'}.  Sealing the outer
        % rows' sides (op:none) removes their own west/east stubs, leaving
        % K_i's west/east open as the only horizontal legs.  Boundary signature
        % (virtual-west=1 | virtual-east=1 | physical-up=1 | physical-down=1);
        % the right-hand side 0 is the zero element of the same matrix space,
        % written as plain text since it carries no ink of its own.
        \tenkzkernel
        \begin{tenkz}[rows={op,op,op}, cols=1, physical=updown]
            \tn[skin=mpo]{P_j} \\
            \tn[skin=mpo,
                ports={180:virtual, 0:virtual}]{\noexpand\mathcal K_i} \\
            \tn[skin=mpo]{P_{j^\prime}}
        \end{tenkz}
        $ = 0\quad\text{unless }i=j=j^\prime$
    \end{center}
    \caption{The local-orthogonality contraction of
    \cite[Appendix~C.2, lines~1457--1470]{Cirac2016MPDO_arXiv}. After
    decomposing the sector trace matrix into primitive components, the source
    assumes that two components occur, obtains two locally orthogonal
    summands, and chooses orthogonal projectors $P_1$ and $P_2$. The
    displayed identity gives the cross contractions
    $P_1{\cal K}P_2=P_2{\cal K}P_1=0$, and the corresponding two-site
    contraction also vanishes. For
    ${\cal A}_i=\spn\{\tr(XP_i{\cal K}P_i):X\}$, these vanishings and
    injectivity imply
    ${\cal A}_1+{\cal A}_2=\MN{D}$ and
    ${\cal A}_1{\cal A}_2=0$, a contradiction.
    Theorem~\ref{thm:mpdo_active_neighboring_support_strongly_connected} does
    not construct $P_1$ or $P_2$: an arbitrary reachability cut
    $I_*=S\sqcup C$ supplies only ${\cal A}_S{\cal A}_C=0$. The common final
    step is the same one-sided obstruction: two nonzero matrix spaces cannot
    sum to $\MN{D}$ while their ordered products vanish.}
    \label{fig:mpdo_local_orthogonality}
\end{figure}

\begin{theorem}[Positive entries of the active trace matrix]%
    \label{thm:mpdo_active_trace_matrix_positive_entries}
    \lean{MPOTensor.PhysicalSectorFactorization.activeSectorTraceMatrix_nonneg}
    \lean{MPOTensor.PhysicalSectorFactorization.activeSectorTraceMatrix_pos_iff}
    \leanok
    \uses{def:mpdo_active_sector_trace_matrix}
    If every neighboring operator is positive semidefinite, then $T^*$ is
    entrywise non-negative and
    $T^*_{k,h}>0\Longleftrightarrow\eta_{k,h}\ne0$.
\end{theorem}

% **Local fix (periodicity):** the source blocks to remove periodic
% components.  The following two-cycle is an auxiliary substitute, to be
% paired with the three-cycle supplied by triangle closure.  See
% docs/paper-gaps/cpgsv17_mpdo_sal_zcl_eta_local_structure.tex.
\begin{theorem}[Two-step return through every active sector]%
    \label{thm:mpdo_active_sector_two_step_return}
    \lean{MPOTensor.PhysicalSectorFactorization.exists_active_two_cycle}
    \lean{MPOTensor.PhysicalSectorFactorization.%
      activeSectorTraceMatrix_pow_two_diag_pos}
    \leanok
    \uses{thm:mpdo_active_one_site_matrix_spanning,
      thm:mpdo_active_trace_matrix_positive_entries}
    Suppose that the active one-site matrices span $\MN{D}$ and each active
    sector contains a nonzero such matrix. For every $k\in I_*$ there is an
    $h\in I_*$ such that
    \begin{align}
        \eta_{k,h}&\ne0,
        \notag\\
        \eta_{h,k}&\ne0.
        \notag
    \end{align}
    If the neighboring operators are positive semidefinite, then
    $(T^{*2})_{k,k}>0$.
\end{theorem}

\begin{theorem}[Irreducibility of the active trace matrix]%
    \label{thm:mpdo_active_trace_matrix_irreducible}
    \lean{MPOTensor.PhysicalSectorFactorization.activeSectorTraceMatrix_isIrreducible}
    \leanok
    \uses{thm:mpdo_active_neighboring_support_strongly_connected,
      thm:mpdo_active_trace_matrix_positive_entries}
    Suppose that the neighboring operators are positive semidefinite, the
    active one-site matrices span $\MN{D}$, and each active sector contains a
    nonzero such matrix. Then the non-negative matrix $T^*$ is irreducible.
\end{theorem}

% **Local fix (periodicity):** the source blocks to remove periodic
% components.  The following length-three return and common-power argument
% replace that step.  See
% docs/paper-gaps/cpgsv17_mpdo_sal_zcl_eta_local_structure.tex.
\begin{theorem}[Three-step returns and a common positive power]%
    \label{thm:mpdo_active_sector_three_step_return_and_primitivity}
    \lean{MPOTensor.PhysicalSectorFactorization.%
      activeSectorTraceMatrix_pow_three_diag_pos}
    \lean{MPOTensor.PhysicalSectorFactorization.%
      activeSectorTraceMatrix_isPrimitive}
    \leanok
    \uses{thm:matrix_primitive_of_two_three_returns,
      thm:mpdo_active_sector_two_step_return,
      thm:mpdo_active_trace_matrix_irreducible,
      thm:mpdo_sector_support_triangle_closure}
    Suppose that the neighboring operators are positive semidefinite, the
    active one-site matrices span $\MN{D}$, each active sector contains a
    nonzero such matrix, and every nonzero edge closes through a third active
    sector. Then $(T^{*3})_{k,k}>0$ for every $k\in I_*$. Together with the
    length-two returns and irreducibility, this implies that $T^*$ is
    primitive.
\end{theorem}

\begin{proof}
    \leanok
    Triangle closure supplies a positive closed walk of length three through
    each active sector. Combine these walks with the length-two returns and
    apply Theorem~\ref{thm:matrix_primitive_of_two_three_returns}.
\end{proof}

% The effective-sector restriction, including the zero rows and columns on
% the full ambient sector set, is documented in
% docs/paper-gaps/cpgsv17_mpdo_sal_zcl_eta_local_structure.tex.
\begin{theorem}[Primitivity of the active sector trace matrix]%
    \label{thm:mpdo_sector_trace_matrix_primitive}
    \lean{MPOTensor.zeroWeightReparameterized_sectorVirtualMatrix_eq_zero}
    \lean{MPOTensor.%
      exists_rephased_inverseMap_activeSectorTraceMatrix_isPrimitive}
    \lean{MPOTensor.%
      exists_positive_physicalSectorFactorization_activeSectorTraceMatrix_isPrimitive_of_isSAL}
    \leanok
    \uses{def:mpdo, def:injective, def:is_sal,
        def:mpdo_active_sector_trace_matrix,
        def:mpdo_explicit_eta_of_sector_tensors,
        thm:mpdo_eta_coherent_positive_choice,
        thm:mpdo_sector_eta_data_trace_matrix}
    Let ${\cal K}$ be an injective tensor that generates MPDOs and
    satisfies SAL, and choose the positive semidefinite neighboring operators of
    Theorem~\ref{thm:mpdo_eta_coherent_positive_choice}. Restrict the sector
    index to the positive-weight summands of the Hayashi decomposition. Then
    positive semidefiniteness makes every neighboring trace real and
    non-negative, and the matrix
    \begin{align}
        T^*_{k,h}
        &=\re\tr(\eta_{k,h})
        \notag
    \end{align}
    is primitive. Zero-weight summands remain in the physical direct sum but
    are not indices of $T^*$.

\end{theorem}

\begin{proof}
    \uses{def:mpdo_eta_recurrent_support,
        def:mpdo_zero_weight_reparameterized_inverse_map_factorization,
        thm:mpdo_zero_weight_reparameterized_neighboring_operator,
        thm:mpdo_sector_virtual_matrix_spanning,
        thm:mpdo_active_one_site_matrix_spanning,
        thm:mpdo_sector_support_triangle_closure,
        thm:mpdo_active_neighboring_support_strongly_connected,
        thm:mpdo_active_sector_two_step_return,
        thm:mpdo_active_trace_matrix_irreducible,
        thm:mpdo_active_trace_matrix_positive_entries,
        thm:mpdo_active_sector_three_step_return_and_primitivity}
    \leanok
    First use the zero-weight reparameterization to make every zero-weight
    sector isolated. Injectivity makes the nonzero-weight one-site matrices
    span $\MN{D}$, and every such sector contains a nonzero matrix. Hence
    Theorem~\ref{thm:mpdo_active_neighboring_support_strongly_connected}
    gives strong connectivity of the active neighboring support. Its proof is
    the one-sided directed-cut argument at
    \cite[Appendix~C.2, lines~1463--1470]{Cirac2016MPDO_arXiv}.

    Finally, nondegeneracy and cyclicity of the trace pairing give a closed
    support walk of length two through every active sector. Triangle closure
    gives one of length three. Strong connectivity together with these two
    coprime return lengths yields a common length at which every entry of the
    active trace-matrix power is positive. Thus $T$ is primitive.
\end{proof}

\begin{theorem}[Zero-correlation-length relations on the primitive active trace matrix]%
    \label{thm:mpdo_active_trace_matrix_zcl_relations}
    \lean{MPOTensor.PhysicalSectorFactorization.%
      activeSectorTraceMatrix_normalized_relations_of_isSourceZCL}
    \lean{MPOTensor.%
      exists_rephased_inverseMap_activeSectorTraceMatrix_normalized_relations}
    \lean{MPOTensor.%
      exists_activeTraceMatrix_relations_of_isSAL_of_isSourceZCL}
    \leanok
    \uses{def:mpdo_active_sector_trace_matrix,
      thm:mpdo_sector_trace_matrix_primitive,
      thm:mpdo_closed_sector_pairing_normalized_idempotent}
    Let ${\cal K}$ be injective and suppose that it satisfies SAL and source
    zero correlation length. Then there are a physical-sector factorization,
    weights $p_k\geq0$ with $\sum_k p_k=1$, and positive semidefinite
    neighboring operators $\eta_{k,h}$ such that the active trace matrix
    \begin{align}
        T^*_{k,h}
        &=\re\tr(\eta_{k,h}),
        \qquad k,h\in I_*:=\{j:p_j\ne0\},
        \notag
    \end{align}
    is primitive. There is also a real number $\lambda>0$ for which, upon
    setting $\widehat T^*:=\lambda^{-1}T^*$, one has
    $(\widehat T^*)^2=(\widehat T^*)^3$. For every positive integer $N$,
    $\tr((\widehat T^*)^N)=\tr(\widehat T^*)$. Thus the primitivity and the
    two normalized identities concern the same active, positively rephased
    trace matrix. No idempotence, rank-one, or semisimplicity conclusion is
    asserted.
\end{theorem}

\begin{proof}
    \leanok
    \uses{thm:mpdo_sector_trace_matrix_primitive,
      thm:mpdo_closed_sector_pairing_normalized_idempotent}
    Choose the coherent sector phases from
    Theorem~\ref{thm:mpdo_sector_trace_matrix_primitive}. For the resulting
    factorization, let $L$ have columns
    $(\tr((l_h)_\beta))_\beta$ and let $Q$ have rows
    $(\tr((r_k)_\alpha))_\alpha$, with both indices restricted to
    $I_*$. The left tensor vanishes when $p_k=0$, so the rectangular product
    $LQ$ is the full physical-trace transfer. Source zero correlation length
    therefore makes $\lambda^{-1}LQ$ idempotent for some $\lambda>0$.

    Positivity of the neighboring operators makes their traces real, and
    hence $QL$ is the complexification of $T^*$. Associativity and cyclicity
    of trace give, for every $N\geq1$,
    \begin{align}
        (\lambda^{-1}QL)^2
        &=(\lambda^{-1}QL)^3,
        \notag\\
        \tr((\lambda^{-1}QL)^N)
        &=\tr(\lambda^{-1}QL).
        \notag
    \end{align}
    Injectivity of the real-to-complex inclusion gives the two real matrix
    identities.
\end{proof}

\begin{theorem}[Virtual spanning does not force active-sector idempotence]
    \label{thm:mpdo_virtual_spanning_active_trace_counterexample}
    \lean{MPOTensor.ActiveSectorSpanningCounterexample.%
      virtual_spanning_does_not_force_rectangular_remainder_zero}
    \lean{MPOTensor.ActiveSectorSpanningCounterexample.%
      physTraceTransfer_tensor}
    \lean{MPOTensor.ActiveSectorSpanningCounterexample.%
      reindex_activeSectorTraceMatrix}
    \leanok
    \uses{def:mpdo_physical_sector_factorization,
      def:mpdo_active_sector_trace_matrix}
    There is a tensor of bond dimension two with four one-dimensional physical
    sectors and weights $p_k$ such that the following statements hold
    simultaneously: the tensor is injective and has source zero correlation
    length; its physical-trace transfer is idempotent; for every
    $k\in\{0,1,2,3\}$, the weights satisfy
    \begin{align}
        p_k&=\frac14,
        \notag\\
        0&\leq p_k,
        \notag\\
        \sum_{k=0}^{3}p_k&=1;
        \notag
    \end{align}
    its sector virtual matrices span $\MN{2}$; every neighboring operator is
    positive semidefinite; and its active trace matrix is primitive and has
    trace one. Nevertheless, if
    $S=LQ$ is the physical-trace transfer and $T=QL$ is the active trace matrix,
    then $Q(1-S)L=T-T^2\ne0$. In particular, full virtual-matrix spanning does
    not by itself imply that the active trace matrix is idempotent.
\end{theorem}

\begin{proof}
    \leanok
    Take
    \begin{align}
        L&=\begin{pmatrix}
            \frac14&\frac14&\frac14&\frac14\\
            1&-1&1&-1
        \end{pmatrix},
        \notag\\
        Q&=\begin{pmatrix}
            1&\frac18\\
            1&\frac18\\
            1&-\frac18\\
            1&-\frac18
        \end{pmatrix}.
        \notag
    \end{align}
    For each $k$, the $k$th sector virtual matrix is the outer product of the
    $k$th column of $L$ with the $k$th row of $Q$. Their four coordinate
    functions are proportional to the Walsh functions $1,x,y,xy$ on four
    points. More explicitly, the matrix whose columns are their coordinates
    in the order $(00,01,10,11)$ has determinant $-1/64$. Hence they form a
    basis of $\MN{2}$. Placing them on the
    diagonal physical entries gives an injective tensor and a physical-sector factorization with
    scalar sectors.

    Direct multiplication gives
    \begin{align}
        LQ
        &=\begin{pmatrix}
            1&0\\
            0&0
        \end{pmatrix},
        \notag\\
        QL
        &=\begin{pmatrix}
            \frac38&\frac18&\frac38&\frac18\\
            \frac38&\frac18&\frac38&\frac18\\
            \frac18&\frac38&\frac18&\frac38\\
            \frac18&\frac38&\frac18&\frac38
        \end{pmatrix}.
        \notag
    \end{align}
    Thus $LQ$ is idempotent, whereas every entry of $QL$ is strictly positive,
    so $QL$ is primitive. Its trace is one. Since each sector is
    one-dimensional, the neighboring operators are the scalar matrices
    $\eta_{k,h}=\begin{pmatrix}(QL)_{k,h}\end{pmatrix}\succeq0$.

    The $(0,1)$ entry of $(QL)^2$ is $1/4$, while the corresponding entry of
    $QL$ is $1/8$. Hence $QL$ is not idempotent, and associativity gives
    $Q(1-LQ)L=QL-(QL)^2\ne0$.
\end{proof}

% Project derivation preparing the Lemma C.5 route of the eta-local
% structure note; no CPSV16 statement is claimed.  See
% docs/paper-gaps/cpgsv17_mpdo_sal_zcl_eta_local_structure.tex.
\begin{theorem}[Doubled-index self-overlap is the Hilbert--Schmidt trace]
    \label{thm:mpdo_doubled_overlap_hilbert_schmidt}
    \lean{MPOTensor.pairCfgEquiv,
      MPOTensor.mpvOverlap_toMPSTensor_self}
    \leanok
    \uses{def:mpv_overlap, def:mpo_to_mps, def:mpo_family}
    Let $K$ be an MPO tensor. For every chain length $N$, the
    self-overlap of the doubled-index MPS family equals the
    Hilbert--Schmidt trace of the closed operator:
    \begin{align}
        O_{K^{\mathrm{MPS}}K^{\mathrm{MPS}}}(N)
        &=\tr\!\left(\rho^{(N)}(K)\,\rho^{(N)}(K)^\dagger\right).
        \notag
    \end{align}
    This is a project derivation preparing the Lemma~C.5 route of
    \cite[Appendix~C.2, lines~1473--1499]{Cirac2016MPDO_arXiv}; the
    pure-state analogue is the normalized-MPV self-overlap lemma of
    \cite[lines~1080--1100]{Cirac2016MPDO_arXiv}. It is not a statement
    of that source. Until a trace normalization is established, the
    right-hand side is the raw squared Hilbert--Schmidt norm of the MPO
    family, not a normalized-state purity.
\end{theorem}

\begin{proof}
    \leanok
    \uses{def:mpv_overlap, def:mpo_to_mps, def:mpo_family}
    Both sides are sums over physical configurations. Expanding the
    trace of the product gives the Frobenius inner product
    \begin{align}
        \tr\!\left(\rho^{(N)}(K)\,\rho^{(N)}(K)^\dagger\right)
        &=\sum_{\sigma,\tau}
          \rho^{(N)}(K)_{\sigma,\tau}\,
          \overline{\rho^{(N)}(K)_{\sigma,\tau}}
        \notag
    \end{align}
    over pairs $(\sigma,\tau)$ of ket and bra configurations, while the
    overlap sums
    $V^{(N)}(K^{\mathrm{MPS}})_\rho\,
    \overline{V^{(N)}(K^{\mathrm{MPS}})_\rho}$ over doubled-index words
    $\rho$. The product encoding of Definition~\ref{def:mpo_to_mps}
    reindexes each configuration pair $(\sigma,\tau)$ as the
    doubled-index word $n\mapsto(\sigma_n,\tau_n)$, and decoding each
    doubled index as a ket--bra pair identifies the summands:
    $V^{(N)}(K^{\mathrm{MPS}})_{(\sigma,\tau)}
    =\rho^{(N)}(K)_{\sigma,\tau}$.
\end{proof}

\begin{corollary}[Self-overlap as the trace of the square for an MPDO]
    \label{cor:mpdo_doubled_overlap_trace_sq}
    \lean{MPOTensor.mpvOverlap_toMPSTensor_self_eq_trace_sq}
    \leanok
    \uses{thm:mpdo_doubled_overlap_hilbert_schmidt, def:mpdo}
    If $K$ generates MPDOs and $N\geq1$, then
    \begin{align}
        O_{K^{\mathrm{MPS}}K^{\mathrm{MPS}}}(N)
        &=\tr\!\left(\rho^{(N)}(K)\,\rho^{(N)}(K)\right).
        \notag
    \end{align}
\end{corollary}

\begin{proof}
    \leanok
    \uses{thm:mpdo_doubled_overlap_hilbert_schmidt, def:mpdo}
    Positivity on the nonempty chain makes $\rho^{(N)}(K)$ positive
    semidefinite, hence Hermitian. Substituting
    $\rho^{(N)}(K)^\dagger=\rho^{(N)}(K)$ into
    Theorem~\ref{thm:mpdo_doubled_overlap_hilbert_schmidt} gives the
    identity.
\end{proof}

\begin{corollary}[Real part of the self-overlap is the squared Frobenius
    norm]
    \label{cor:mpdo_doubled_overlap_frobenius}
    \lean{MPOTensor.mpvOverlap_toMPSTensor_self_re_eq_frobenius_norm_sq}
    \leanok
    \uses{thm:mpdo_doubled_overlap_hilbert_schmidt}
    For every MPO tensor $K$ and every chain length $N$,
    \begin{align}
        \re O_{K^{\mathrm{MPS}}K^{\mathrm{MPS}}}(N)
        &=\|\rho^{(N)}(K)\|_F^2
          =\sum_{\sigma,\tau}
          \left|\rho^{(N)}(K)_{\sigma,\tau}\right|^2.
        \notag
    \end{align}
\end{corollary}

\begin{proof}
    \leanok
    \uses{thm:mpdo_doubled_overlap_hilbert_schmidt}
    Cyclicity of the trace rewrites the Hilbert--Schmidt trace as
    $\tr\!\left(\rho^{(N)}(K)^\dagger\rho^{(N)}(K)\right)$, whose real
    part is the sum of the squared absolute entries.
\end{proof}

\begin{theorem}[Trace factorization of the cyclic neighboring product]
    \label{thm:mpdo_cyclic_neighboring_trace_products}
    \lean{MPOTensor.PhysicalSectorFactorization.%
      trace_cyclicNeighboringProduct_eq_prod_trace,
      MPOTensor.PhysicalSectorFactorization.%
      trace_cyclicNeighboringProduct_sq_eq_prod_trace_sq}
    \leanok
    \uses{def:mpdo_physical_sector_chain_coordinates}
    Let $N\geq1$ and let $k=(k_0,\ldots,k_{N-1})$ be a cyclic sector
    configuration of a physical-sector factorization. The trace of the
    cyclic neighboring product factors into the product of the
    neighboring traces, and so does the trace of its square:
    \begin{align}
        \tr(\Omega_k)
        &=\prod_{n=0}^{N-1}\tr\!\left(\eta_{k_n,k_{n+1}}\right),
        \notag\\
        \tr\!\left(\Omega_k^2\right)
        &=\prod_{n=0}^{N-1}
          \tr\!\left(\eta_{k_n,k_{n+1}}^2\right),
        \notag
    \end{align}
    with site labels read modulo $N$. These are the
    coefficient-extraction identities for the cyclic direct-sum form of
    \cite[Appendix~C.2, lines~1580--1593]{Cirac2016MPDO_arXiv}.
\end{theorem}

\begin{proof}
    \leanok
    \uses{def:mpdo_physical_sector_chain_coordinates}
    The fiber bijection of
    Definition~\ref{def:mpdo_physical_sector_chain_coordinates}
    identifies the diagonal entries of $\Omega_k$ with tuples of
    diagonal entries of the neighboring factors. Hence
    \begin{align}
        \tr(\Omega_k)
        &=\sum_x\prod_{n=0}^{N-1}
          \eta_{k_n,k_{n+1}}(x_n,x_n)
          =\prod_{n=0}^{N-1}\sum_{x_n}
          \eta_{k_n,k_{n+1}}(x_n,x_n),
        \notag
    \end{align}
    expanding the product of sums. For the square, write
    $\tr(\Omega_k^2)=\sum_{x,y}\Omega_k(x,y)\,\Omega_k(y,x)$; the same
    reindexing and expansion turn the inner sums into the diagonal
    entries of $\eta_{k_n,k_{n+1}}^2$.
\end{proof}

% Proved Case-I components; the singleton-sector assembly (the Case-I
% relations $(T^*)^2=T^*$ and $Q(1-LQ)L=0$) is not yet formalized.  The
% route is sketched in docs/paper-gaps/cpsv16_commuting_form_to_sal.tex.
\begin{lemma}[Trace-squared matrix of the neighboring operators]
    \label{lem:mpdo_active_sector_trace_sq_matrix}
    \lean{MPOTensor.activeSectorTraceSqMatrix,
      MPOTensor.activeSectorTraceSqMatrix_nonneg,
      MPOTensor.activeSectorTraceSqMatrix_le_activeSectorTraceMatrix_sq}
    \leanok
    \uses{def:mpdo_active_sector_trace_matrix}
    Let ${\cal K}$ be an MPO tensor with a physical-sector factorization
    and sector weights $p_k$, and suppose that every neighboring
    operator $\eta_{k,h}$ is positive semidefinite. On the active
    sector set $I_*$, define the trace-squared matrix
    \begin{align}
        S^*_{k,h}
        &:=\re\tr\!\left(\eta_{k,h}^2\right),
        \qquad k,h\in I_*.
        \notag
    \end{align}
    Then, entrywise,
    \begin{align}
        0&\leq S^*_{k,h}\leq\left(T^*_{k,h}\right)^2.
        \notag
    \end{align}
    This matrix is a project derivation preparing the Case-I route
    toward
    \cite[Appendix~C.2, Lemma~C.5, lines~1473--1499]{Cirac2016MPDO_arXiv};
    it is not a matrix defined in that source.
\end{lemma}

\begin{proof}
    \leanok
    \uses{def:mpdo_active_sector_trace_matrix,
      thm:psd_trace_square_le_square_trace}
    Hermiticity of $\eta_{k,h}$ gives
    $\eta_{k,h}^2=\eta_{k,h}\,\eta_{k,h}^\dagger$, which is positive
    semidefinite, so its trace is a non-negative real number. The upper
    bound is Theorem~\ref{thm:psd_trace_square_le_square_trace} applied
    to $\eta_{k,h}$.
\end{proof}

% Proved Case-I components.  The remaining singleton-sector assembly
% is not yet formalized: the Case-I relations $(T^*)^2=T^*$ and
% $Q(1-LQ)L=0$ in the rectangular decomposition
% $\mathcal T_{\cal K}=LQ$.  The virtual-spanning counterexample above
% shows that the last relation does not follow from virtual spanning
% alone.  The route is sketched in
% docs/paper-gaps/cpsv16_commuting_form_to_sal.tex.
\begin{theorem}[Literal zero correlation length stabilizes the active
    trace matrix]
    \label{thm:mpdo_active_trace_matrix_literal_zcl_powers}
    \lean{MPOTensor.%
      activeSectorTraceMatrix_pow_two_eq_pow_three_of_literal_ZCL,
      MPOTensor.activeSectorTraceMatrix_pow_two_pos}
    \leanok
    \uses{def:mpdo_active_sector_trace_matrix,
      def:mpo_physical_trace_transfer}
    Let ${\cal K}$ be an MPO tensor with a physical-sector factorization
    and sector weights $p_k$. Suppose that every neighboring operator
    is positive semidefinite, that the physical-trace transfer is
    literally idempotent,
    \begin{align}
        \mathcal T_{\cal K}^2&=\mathcal T_{\cal K},
        \notag
    \end{align}
    as in \cite[Definition~4.2, lines~735--741]{Cirac2016MPDO_arXiv},
    and that the left tensor of every zero-weight sector vanishes. Then
    the active trace matrix stabilizes after the second power:
    \begin{align}
        \left(T^*\right)^2&=\left(T^*\right)^3.
        \notag
    \end{align}
    If, in addition, the active one-site matrices span $\MN{D}$, each
    active sector contains a nonzero such matrix, every nonzero
    neighboring operator closes through a third active sector, and
    there is at least one active sector, then $(T^*)^2$ is strictly
    positive entrywise.

    These are matrix-algebra components of the Case-I argument under
    the assumptions of
    \cite[Appendix~C.2, lines~1374--1381]{Cirac2016MPDO_arXiv}, for the
    factorization supplied by
    \cite[Appendix~C.2, Lemma~C.4, lines~1406--1471]{Cirac2016MPDO_arXiv}
    and the zero-correlation-length step of
    \cite[Appendix~C.2, Lemma~C.5, lines~1473--1499]{Cirac2016MPDO_arXiv}.
    They are project derivations, not statements of that source.
\end{theorem}

\begin{proof}
    \leanok
    \uses{def:mpdo_active_sector_trace_matrix,
      thm:mpdo_active_sector_three_step_return_and_primitivity,
      thm:matrix_primitive_pow_two_rank_one}
    On the active sectors, collect the left and right sector-tensor
    traces into rectangular matrices $L$ and $Q$ with
    \begin{align}
        L_{\beta,h}&=\tr\!\left((l_h)_\beta\right),
        &
        Q_{k,\alpha}&=\tr\!\left((r_k)_\alpha\right).
        \notag
    \end{align}
    The left tensor vanishes on every zero-weight sector, so the
    rectangular product $LQ$ is the full physical-trace transfer, and
    literal zero correlation length makes $LQ$ idempotent.
    Associativity then gives
    \begin{align}
        (QL)^2&=Q(LQ)L=Q(LQ)(LQ)L=(QL)^3.
        \notag
    \end{align}
    Positivity of the neighboring operators makes their traces real,
    and expanding the products identifies $QL$ with the
    complexification of $T^*$; injectivity of the real-to-complex
    inclusion gives $(T^*)^2=(T^*)^3$.

    Under the additional hypotheses,
    Theorem~\ref{thm:mpdo_active_sector_three_step_return_and_primitivity}
    makes $T^*$ primitive: some power $m\geq1$ has strictly positive
    entries. Stabilization gives $(T^*)^{2m}=(T^*)^2$, so
    \begin{align}
        \left((T^*)^2\right)_{k,h}
        &=\sum_j\left((T^*)^m\right)_{k,j}
          \left((T^*)^m\right)_{j,h}>0,
        \notag
    \end{align}
    which is the positivity argument of
    Theorem~\ref{thm:matrix_primitive_pow_two_rank_one}.
\end{proof}

% Elementary combinatorial infrastructure for the overlap formula
% $O_{{\cal K}^{\mathrm{MPS}}{\cal K}^{\mathrm{MPS}}}(N)=\tr((S^*)^N)$,
% which is not yet assembled: reconciling the active-sector-restricted
% index set $I_*$ used by $S^*$ with the full sector index set that the
% cyclic-product trace factorization of
% Theorem~\ref{thm:mpdo_cyclic_neighboring_trace_products} ranges over
% remains open.  The route is sketched in
% docs/paper-gaps/cpsv16_commuting_form_to_sal.tex.
\begin{lemma}[Indicator-path expansion of a matrix power]
    \label{lem:mpdo_pow_apply_path_indicator}
    \lean{MPOTensor.pow_apply_eq_sum_path_indicator}
    \leanok
    Let $S$ be a square matrix with entries in a commutative semiring,
    indexed by a finite set $I$, and let $N\geq0$. Then, for $a,b\in I$,
    \begin{align}
        \left(S^N\right)_{a,b}
        &=\sum_{\substack{v:\{0,\ldots,N\}\to I\\ v_0=a,\ v_N=b}}
          \ \prod_{i=0}^{N-1}S_{v_i,v_{i+1}},
        \notag
    \end{align}
    the sum, over length-$N$ walks $v$ in $I$ from $a$ to $b$, of the
    product of matrix entries along the walk. This is an elementary
    combinatorial identity used to derive the cyclic-product form of
    $\tr(S^N)$ in the Case-I route.
\end{lemma}

\begin{proof}
    \leanok
    Induction on $N$: the base case is the identity-matrix entry
    $\delta_{a,b}$, indexed by the unique length-$0$ walk; the
    inductive step expands $\left(S^{N+1}\right)_{a,b}
    =\sum_c\left(S^N\right)_{a,c}S_{c,b}$ and reindexes the
    walk-and-endpoint pair $(v,c)$ as the length-$(N+1)$ walk obtained
    by appending $c$ to $v$.
\end{proof}

\begin{theorem}[Trace of a matrix power as a sum over cyclic products]
    \label{thm:mpdo_trace_pow_cyclic_product}
    \lean{MPOTensor.trace_pow_eq_sum_cyclic_product}
    \leanok
    \uses{lem:mpdo_pow_apply_path_indicator}
    Let $S$ be a square matrix with entries in a commutative semiring,
    indexed by a finite set $I$, and let $N\geq1$. Then
    \begin{align}
        \tr\!\left(S^N\right)
        &=\sum_{k:\{0,\ldots,N-1\}\to I}
          \ \prod_{i=0}^{N-1}S_{k_i,k_{i+1}},
        \notag
    \end{align}
    with indices read modulo $N$, the sum over labelings $k$ of the
    $N$ cyclic positions of the product of matrix entries along the
    cycle. Applied to the trace-squared matrix $S^*$ of
    Lemma~\ref{lem:mpdo_active_sector_trace_sq_matrix}, this is the
    combinatorial identity used to express the self-overlap
    $O_{{\cal K}^{\mathrm{MPS}}{\cal K}^{\mathrm{MPS}}}(N)$ in terms of
    $\tr\!\left((S^*)^N\right)$.
\end{theorem}

\begin{proof}
    \leanok
    \uses{lem:mpdo_pow_apply_path_indicator}
    Expand $\tr(S^N)=\sum_a\left(S^N\right)_{a,a}$ by
    Lemma~\ref{lem:mpdo_pow_apply_path_indicator}, collapse the sum
    over the shared endpoint $a=v_0=v_N$, and reindex the resulting
    closed length-$N$ walks $v$ (with $v_0=v_N$) as cyclic labelings
    $k:\{0,\ldots,N-1\}\to I$ by dropping the repeated final
    coordinate.
\end{proof}

\begin{lemma}[Trace invariance under the physical basis change]
    \label{lem:mpdo_trace_conjugate_physical_invariant}
    \lean{MPOTensor.sitewise_prod_conjTranspose_mul_self,
      MPOTensor.trace_mpo_conjugatePhysical_mul_self}
    \leanok
    \uses{def:mpo_family}
    Let $K$ be an MPO tensor and $U$ an isometry on the physical space,
    $U^\dagger U=1$. Then
    \begin{align}
        \tr\!\left(\rho^{(N)}(\widetilde K)\,\rho^{(N)}(\widetilde K)\right)
        &=\tr\!\left(\rho^{(N)}(K)\,\rho^{(N)}(K)\right),
        \notag
    \end{align}
    where $\widetilde K$ is $K$ conjugated by $U$ at every physical site.
    This is a project derivation preparing the Case-I route of
    \cite[Appendix~C.2, lines~1434--1448]{Cirac2016MPDO_arXiv} (the basis
    change defining $\sigma_{\mathrm{tilde}}$); it is not a claim from that
    source.
\end{lemma}

\begin{proof}
    \leanok
    The sitewise product matrix $W_{\sigma,\sigma'}=\prod_n U(\sigma_n,\sigma'_n)$
    satisfies $W^\dagger W=1$ because $U$ does: expanding $W^\dagger W$ as a sum
    over intermediate configurations and using $\sum_x\overline{U(x,a)}\,U(x,b)
    =\delta_{a,b}$ site by site gives the identity. Since
    $\rho^{(N)}(\widetilde K)=W\,\rho^{(N)}(K)\,W^\dagger$, cyclicity of the
    trace and $W^\dagger W=1$ cancel the $W$'s from
    $\tr\!\left(W\rho^{(N)}(K)\rho^{(N)}(K)W^\dagger\right)$.
\end{proof}

\begin{lemma}[Reduction to the active sector]
    \label{lem:mpdo_sum_prod_traceSq_active}
    \lean{MPOTensor.neighboringOperator_eq_zero_of_inactive_right,
      MPOTensor.prod_traceSq_eq_zero_of_not_forall_active,
      MPOTensor.sum_prod_traceSq_eq_sum_active}
    \leanok
    \uses{def:mpdo_active_sector_trace_matrix,
      thm:mpdo_cyclic_neighboring_trace_products}
    Let ${\cal K}$ be an MPO tensor with a physical-sector factorization
    and sector weights $p_k$, such that the left tensor of every
    zero-weight sector vanishes. Then, for $N\geq1$, the sum over all
    sector labelings of the cyclic trace-square product reduces to the
    sum over active-sector labelings:
    \begin{align}
        \sum_{k:\{0,\ldots,N-1\}\to I}\ \prod_{i=0}^{N-1}
          \tr\!\left(\eta_{k_i,k_{i+1}}^2\right)
        &=\sum_{k:\{0,\ldots,N-1\}\to I_*}\ \prod_{i=0}^{N-1}
          \tr\!\left(\eta_{k_i,k_{i+1}}^2\right).
        \notag
    \end{align}
    This is the reconciliation, flagged as an open reconciliation gap when
    Lemma~\ref{lem:mpdo_pow_apply_path_indicator} and
    Theorem~\ref{thm:mpdo_trace_pow_cyclic_product} were added, between
    the full sector index set $I$ and the active sector set $I_*$.
\end{lemma}

\begin{proof}
    \leanok
    \uses{def:mpdo_active_sector_trace_matrix}
    A neighboring operator with an inactive target sector vanishes,
    because the left tensor of that sector vanishes identically. Hence any
    cyclic labeling touching an inactive sector has a vanishing factor
    (that sector is the target of some edge in the cycle), so its product
    is zero. The active-sector sum is therefore exactly the full-sector
    sum, restricted to a labeling-by-labeling bijection with the inactive
    terms discarded.
\end{proof}

\begin{theorem}[The overlap formula]
    \label{thm:mpdo_overlap_formula}
    \lean{MPOTensor.mpvOverlap_toMPSTensor_self_eq_ofReal_trace_activeSectorTraceSqMatrix_pow}
    \leanok
    \uses{lem:mpdo_active_sector_trace_sq_matrix,
      thm:mpdo_doubled_overlap_hilbert_schmidt,
      thm:mpdo_trace_pow_cyclic_product,
      lem:mpdo_trace_conjugate_physical_invariant,
      lem:mpdo_sum_prod_traceSq_active}
    Let ${\cal K}$ be an MPO tensor with a physical-sector factorization
    and sector weights $p_k$, such that every neighboring operator is
    positive semidefinite and the left tensor of every zero-weight sector
    vanishes. Then, for $N\geq1$,
    \begin{align}
        O_{{\cal K}^{\mathrm{MPS}}{\cal K}^{\mathrm{MPS}}}(N)
        &=\tr\!\left((S^*)^N\right).
        \notag
    \end{align}
    This is a project derivation completing the first step of the Case-I
    route toward
    \cite[Appendix~C.2, Lemma~C.5, lines~1473--1499]{Cirac2016MPDO_arXiv};
    it is not a statement of that source.
\end{theorem}

\begin{proof}
    \leanok
    \uses{lem:mpdo_active_sector_trace_sq_matrix,
      thm:mpdo_doubled_overlap_hilbert_schmidt,
      thm:mpdo_trace_pow_cyclic_product,
      thm:mpdo_cyclic_neighboring_trace_products,
      lem:mpdo_trace_conjugate_physical_invariant,
      lem:mpdo_sum_prod_traceSq_active}
    Write $\rho=\rho^{(N)}({\cal K})$, $\widetilde\rho=\rho^{(N)}(\widetilde{\cal K})$
    for the physically conjugated MPO, and $\Omega_k$ for the cyclic
    neighboring product of Theorem~\ref{thm:mpdo_cyclic_neighboring_trace_products}.
    The overlap unwinds into a chain of trace identities:
    \begin{align}
        O_{{\cal K}^{\mathrm{MPS}}{\cal K}^{\mathrm{MPS}}}(N)
        &=\tr(\rho\,\rho)
        &&\text{Theorem~\ref{thm:mpdo_doubled_overlap_hilbert_schmidt}}
        \notag\\
        &=\tr(\widetilde\rho\,\widetilde\rho)
        &&\text{Lemma~\ref{lem:mpdo_trace_conjugate_physical_invariant}}
        \notag\\
        &=\tr\!\left(\Bigl(\bigoplus_{k}\Omega_k\Bigr)^{\!2}\right)
        &&\text{block-diagonal product realization}
        \notag\\
        &=\sum_{k:\{0,\ldots,N-1\}\to I}\tr\!\left(\Omega_k^2\right)
        &&\text{trace of a block-diagonal matrix}
        \notag\\
        &=\sum_{k:\{0,\ldots,N-1\}\to I}\ \prod_{i=0}^{N-1}
          \tr\!\left(\eta_{k_i,k_{i+1}}^2\right)
        &&\text{Theorem~\ref{thm:mpdo_cyclic_neighboring_trace_products}}
        \notag\\
        &=\sum_{k:\{0,\ldots,N-1\}\to I_*}\ \prod_{i=0}^{N-1}
          \tr\!\left(\eta_{k_i,k_{i+1}}^2\right)
        &&\text{Lemma~\ref{lem:mpdo_sum_prod_traceSq_active}}
        \notag\\
        &=\sum_{k:\{0,\ldots,N-1\}\to I_*}\ \prod_{i=0}^{N-1}S^*_{k_i,k_{i+1}}
        &&\text{Hermiticity, }\tr(A)=\re\tr(A)\text{ for }A=A^\dagger
        \notag\\
        &=\tr\!\left((S^*)^N\right)
        &&\text{Theorem~\ref{thm:mpdo_trace_pow_cyclic_product}}
        \notag
    \end{align}
    The third step combines the block-diagonal product realization of
    the sector-coordinate MPO with the reindexing invariance of
    trace-of-square. The sixth step uses Hermiticity of $\eta_{k,h}$ to
    identify its complex trace-square with the real entry
    $S^*_{k,h}=\re\tr(\eta_{k,h}^2)$.
\end{proof}
